\begin{apartats}

\apartat[1,25]{0,75}
Calcula les integrals $\displaystyle\int(3x^2+4x-2)\,dx$ i $\displaystyle\int(x-1)^3\,dx$.

\begin{solucio}
$\displaystyle\int(3x^2+4x-2)\,dx=x^3+2x^2-2x+k$ i $\displaystyle\int(x-1)^3\,dx=\frac{(x-1)^4}{4}+k$.
\end{solucio}

\begin{tria}{immediates-potencies}
\itemtria{radicals}{1}{1,25}
Calcula $\displaystyle\int\left(2\sqrt x-\frac{1}{\sqrt x}\right)dx$ i
$\displaystyle\int\frac{1}{\sqrt[3]x}\,dx$.

\begin{solucio}
Escrivint les arrels com a potències, $2x^{1/2}-x^{-1/2}$ i $x^{-1/3}$:
\[
\int\left(2x^{1/2}-x^{-1/2}\right)dx=2\,\frac{x^{3/2}}{3/2}-\frac{x^{1/2}}{1/2}+k=\frac43x\sqrt x-2\sqrt x+k,
\qquad
\int x^{-1/3}\,dx=\frac{x^{2/3}}{2/3}+k=\frac32\sqrt[3]{x^2}+k.
\]
\end{solucio}

\itemtria{racionals}{1}{1,25}
Calcula $\displaystyle\int\left(\frac1x+\frac2{x^2}\right)dx$ i
$\displaystyle\int\left(\frac{3}{x+2}-\frac{1}{(x-3)^2}\right)dx$.

\begin{solucio}
$\frac1x$ té primitiva $\ln|x|$, i $\frac2{x^2}=2x^{-2}$ té primitiva $2\frac{x^{-1}}{-1}=-\frac2x$:
$\displaystyle\int\left(\frac1x+\frac2{x^2}\right)dx=\ln|x|-\frac2x+k$.\\
De la mateixa manera, amb $x+2$ i $x-3$ al lloc de $x$:
$\displaystyle\int\left(\frac{3}{x+2}-\frac{1}{(x-3)^2}\right)dx=3\ln|x+2|+\frac{1}{x-3}+k$.
\end{solucio}
\end{tria}

\begin{nomesllarg}
\apartat{0,75}
Calcula $\displaystyle\int(1-2x)^2\,dx$ desenvolupant el quadrat. Un company ha obtingut
$-\frac{(1-2x)^3}{6}+k$. Té raó?

\begin{solucio}
$\displaystyle\int(1-4x+4x^2)\,dx=x-2x^2+\frac43x^3+k$. El company també té raó: la derivada de
$-\frac{(1-2x)^3}{6}$ és $-\frac{3(1-2x)^2\cdot(-2)}{6}=(1-2x)^2$. Les dues primitives es diferencien en
una constant: $-\frac{(1-2x)^3}{6}=-\frac16+x-2x^2+\frac43x^3$.
\end{solucio}
\end{nomesllarg}

\end{apartats}
